\begin{textsample}{POS Dim 2 – rj – Score 24.00 – rj/rj\_em\_1.txt}  \label{ex:f2_pos_231}
APRESENTAÇÃO A Secretaria de Estado de Educação do Rio de Janeiro apresenta o \textbf{Currículo} Referencial do Ensino Médio ( Regular e Educação de Jovens e Adultos ) fundamentado nas premissas da Política Pública Educacional do Novo Ensino Médio, iniciando sua implementação no ano de 2022.
A Lei Federal nº 13.415 de 2017, que \textbf{instituiu} a \textbf{Reforma} do Ensino Médio, em âmbito nacional, alterou a Lei de \textbf{Diretrizes} e Bases da Educação Nacional e estabeleceu uma mudança na estrutura do Ensino Médio, ampliando o tempo mínimo do estudante na escola de 800 horas para 1.000 horas anuais, definindo uma nova organização curricular, mais \textbf{flexível}, que contemple uma Base Nacional Comum Curricular ( BNCC ).
A nova organização curricular possibilita a \textbf{oferta} de diferentes possibilidades de escolhas aos jovens pelos \textbf{Itinerários} \textbf{Formativos}, com foco nas áreas de conhecimento e na formação \textbf{técnica} e \textbf{profissional}.
Em 2018, o Governo Federal \textbf{instituiu} a Base Nacional Comum Curricular do Ensino Médio ( BNCC-EM ) e orientou os sistemas de ensino, as instituições e \textbf{redes} escolares na sua execução em \textbf{regime} de colaboração nos termos do Art. 211 da Constituição Federal de 1988 e Art. 8º da Lei nº 9.394 de 1996 ( Lei de \textbf{Diretrizes} e Bases da Educação - LDB ).
Essa iniciativa auxiliou os estados, os \textbf{municípios} e o Distrito Federal na elaboração de seus \textbf{Currículos} \textbf{alinhados} à BNCC.
O Programa de Apoio à Implementação da Base Nacional Comum Curricular ( ProBNCC ) foi \textbf{instituído} por meio da \textbf{Portaria} MEC nº 331, de 5 de abril de 2018, para o apoio \textbf{técnico} e financeiro às Secretarias \textbf{Estaduais} e Distrital de Educação – SEEDUC e às Secretarias \textbf{Municipais} de Educação, no processo de revisão, elaboração e \textbf{implantação} de seus \textbf{Currículos} \textbf{alinhados} à BNCC.
Participaram desse processo o Conselho Nacional de \textbf{Secretários} de Educação ( Consed ) e a União Nacional dos \textbf{Dirigentes} \textbf{Municipais} de Educação ( Undime ), que integram o Comitê Nacional de Implementação da BNCC, as instituições e organizações da Sociedade Civil, as representações institucionais dos Conselhos Nacionais, \textbf{Estaduais} e \textbf{Municipais} ( Conselho Nacional de Educação - CNE, Fórum Nacional dos Conselhos \textbf{Estaduais} de Educação – FNCE e União Nacional dos Conselhos \textbf{Municipais} de Educação - UNCME)¹.
No ano de 2019, a Secretaria Estadual de Educação do Rio de Janeiro – SEEDUC-RJ iniciou a elaboração do Documento de Orientação Curricular do Rio de Janeiro ( DOC-RJ – Etapa Ensino Médio ).
Em janeiro de 2021, o DOC-RJ foi encaminhado para apreciação do Conselho Estadual de Educação do Rio de Janeiro ( CEE-RJ ).
Em 7 de dezembro de 2021, por meio da \textbf{Deliberação} nº 394, o Conselho Estadual de Educação – CEE-RJ \textbf{instituiu} as \textbf{Diretrizes} para a \textbf{implantação} do Documento de Orientação Curricular do Estado do Rio de Janeiro – Ensino médio ( DOC-RJ ) e definiu princípios e referências curriculares para as instituições de educação básica que integram o sistema \textbf{estadual} de ensino do Rio de Janeiro.
A partir dessa \textbf{deliberação}, a SEEDUC-RJ organizou o \textbf{Currículo} Referencial do Estado do Rio de Janeiro – Ensino Médio, disponibilizando -o à apreciação dos docentes da \textbf{rede}, por meio de uma consulta pública, para em seguida, efetivá -lo no ano de 2022.

% matched lemmas: alinhar, currículo, deliberação, diretriz, dirigente, estadual, flexível, formativo, implantação, instituir, itinerário, municipal, município, oferta, portaria, profissional, rede, reforma, regime, secretário, técnico
\end{textsample}
